\documentclass[12pt,a4paper]{article}

\usepackage[hidelayout]{../../../Latex/Style/coursework-report-core}

\addbibresource{res-ml.bib}

\title{RES-ML --- Machine-Learning Methodology}
\author{}
\date{\today}

\begin{document}

\maketitle

\tableofcontents
\clearpage

\section{Domain Purpose and Scope}
\label{sec:res-ml-domain-purpose}

RES-ML defines the machine-learning methodology for the revised research
scope. Under the 2026-07-19 scope decision, ML is deferred future
work. The active Studentship objective is to stabilise the 2011
Tohoku hydrodynamic pipeline, Kamaishi/Sendai corridor comparison and
fixed-rigid barrier-class metrics before any supervised-learning,
surrogate optimisation or generative geometry work begins.

The high-fidelity simulation workflow remains the authoritative source of
physical response. A future surrogate may be an assessment, screening
or sampling tool; it shall not be represented as replacing validated
fluid simulation, structural analysis or direct confirmation.

Machine-learning training is blocked until the selected tsunami event and the
relevant hydrodynamic models have passed the V1--V2--V3 validation
hierarchy recorded in the current scope decision. Historical or
reconstructed defence geometries are evidence about plausible cases,
not direct causal labels for optimal reinforcement; controlled,
validated simulations remain the only acceptable future training
targets.

\subsection{Active Scope}
\label{sec:res-ml-active-scope}

There is no active ML training scope in the current Research phase.
Future RES-ML work may cover:
\begin{itemize}
  \item response surrogates for validated hydrodynamic outputs;
  \item optional adequacy or damage classification after structural
    extensions exist;
  \item screening within a bounded barrier-class or intervention space;
  \item active selection of additional validated simulations;
  \item dimensionality reduction where correlated outputs or histories justify
    it;
  \item small-data algorithm assessment and non-ML baselines;
  \item predictive uncertainty and out-of-domain detection;
  \item direct-simulation fallback for uncertain or unsupported queries;
  \item integration with validated regional and local CFD workflows,
    and later structural/FSI workflows if activated.
\end{itemize}

\subsection{Deferred Scope}
\label{sec:res-ml-deferred-scope}

The following are deferred to later research and shall not be represented as
initial acceptance requirements:
\begin{itemize}
  \item unrestricted free-form geometry generation;
  \item surrogate optimisation during the active hydrodynamic phase;
  \item reinforcement learning;
  \item unrestricted inverse design;
  \item claims of cross-event generalisation;
  \item direct training from before/after geometry alone;
  \item replacement of high-fidelity confirmation.
\end{itemize}

\subsection{Supervised-Learning Relationship}
\label{sec:res-ml-supervised-learning-relationship}

The intended supervised-learning relationship is:
\begin{align}
  &\text{corridor descriptors}
  + \text{fixed-rigid geometry descriptors}
  + \text{validated hydrodynamic-condition descriptors}
  \notag\\
  &\qquad \rightarrow
  \text{hydrodynamic response}
  + \text{predictive uncertainty}
  \label{eq:res-ml-supervised-learning-relationship}
\end{align}

Decision role: \cref{eq:res-ml-supervised-learning-relationship} is a
future data relationship, not an active implementation requirement.
Structural response and damage labels are excluded until RES-STR
activates and validates those extensions.

\subsection{Research Gates}
\label{sec:res-ml-research-gates}

The ML workstream is governed by the following gates:
\begin{enumerate}
  \item \textbf{G1 --- Data-source feasibility.} Source data,
    CRS/datum, licence and preprocessing provenance are recorded. No
    ML training is permitted before this gate.
  \item \textbf{G2 --- Corridor definition.} Kamaishi and Sendai
    corridor geometry is locked. ML remains deferred.
  \item \textbf{G3 --- Regional replay readiness.} The Tohoku regional
    model and replay-boundary histories are ready for validation. No
    ML dataset may be accepted before this gate.
  \item \textbf{G4 --- Local hydrodynamic metric extraction.}
    Fixed-rigid local runs emit complete quality and metric labels.
    Future dataset design may begin, but model training remains
    blocked until validation status is known.
  \item \textbf{G5 --- V1--V2--V3 validation.} Solver capability,
    Tohoku hydrodynamics and local impact/loading evidence are
    accepted. Future ML training may be reconsidered after this gate.
  \item \textbf{G6 --- Barrier-class demonstration.} Barrier classes
    are compared with complete provenance. This is not surrogate
    validation.
\end{enumerate}

\section{Domain Deliverables}
\label{sec:res-ml-domain-deliverables}

% -------------------------------------------------------------------------
% WBS ID: RES-ML-ROL
% Deliverable: Machine-Learning Role and Use-Case Definition
% Status: In progress
% Evidence status: ML deferred by current hydrodynamic scope
% Review status: Not reviewed
% Last updated: 2026-07-19
% Dependencies: RES-DAT-CAS, RES-GEO-SPEC, RES-VER-SPEC
% Downstream interfaces: Future RES-ML-DAT, RES-ML-SUR, RES-ML-SPEC
% -------------------------------------------------------------------------

\section{RES-ML-ROL --- Machine-Learning Role and Use-Case Definition}
\label{sec:res-ml-rol}

\subsection{Purpose and Scope}
\label{sec:res-ml-rol-purpose}

This Deliverable records ML as a deferred future role under the
current Tohoku hydrodynamic scope. The active project shall first
produce validated fixed-rigid barrier hydrodynamic outputs for the
Kamaishi and Sendai corridors.

Deferred possible use cases are:
\begin{itemize}
  \item a forward response surrogate;
  \item a future damage or adequacy classifier after structural
    extensions exist;
  \item barrier-class or intervention screening and ranking;
  \item active-learning support for selecting additional direct simulations;
  \item reduced-order representation of correlated response outputs.
\end{itemize}

Deferred use cases are:
\begin{itemize}
  \item generative geometry;
  \item unrestricted inverse design;
  \item reinforcement learning;
  \item cross-event transfer learning.
\end{itemize}

\subsection{Research Questions}
\label{sec:res-ml-rol-research-questions}

% TODO: State the questions that must be answered before the Deliverable
% can be accepted.

\subsection{Terminology and Definitions}
\label{sec:res-ml-rol-definitions}

% TODO: Define the technical terminology, symbols, quantities and
% classifications used in this Deliverable.

\subsection{Literature Evidence}
\label{sec:res-ml-rol-literature-evidence}

% TODO: Record evidence from peer-reviewed papers, authoritative datasets,
% standards, technical reports and primary documentation.
%
% Do not create a detached literature-summary list. Explain how each source
% supports, contradicts or limits a project decision.

\subsection{Governing Relations and Quantitative Definitions}
\label{sec:res-ml-rol-governing-relations}

% TODO: Record relevant equations, dimensionless groups, empirical models,
% extraction rules or quantitative definitions.
%
% Use align environments for displayed equations.
% Every displayed equation must have a unique \label{eq:...}.
% Do not place punctuation after displayed equations.

\subsection{Machine-Learning Role}
\label{sec:res-ml-rol-machine-learning-role}

Each retained use case shall record:
\begin{itemize}
  \item input features;
  \item prediction or decision target;
  \item required dataset and simulation provenance;
  \item non-ML baseline;
  \item expected advantage over the baseline;
  \item consequence of an incorrect prediction;
  \item uncertainty requirement;
  \item direct-simulation fallback.
\end{itemize}

The consequence analysis shall distinguish high-consequence unsafe predictions
from conservative but inefficient safe-to-unsafe predictions.

\subsection{Non-ML Baseline}
\label{sec:res-ml-rol-non-ml-baseline}

% TODO: Complete this RES-ML Domain-specific subsection.

\subsection{Input and Output Representation}
\label{sec:res-ml-rol-input-output-representation}

% TODO: Complete this RES-ML Domain-specific subsection.

\subsection{Dataset Requirements}
\label{sec:res-ml-rol-dataset-requirements}

% TODO: Complete this RES-ML Domain-specific subsection.

\subsection{Model or Architecture Candidates}
\label{sec:res-ml-rol-model-architecture-candidates}

% TODO: Complete this RES-ML Domain-specific subsection.

\subsection{Training and Validation Method}
\label{sec:res-ml-rol-training-validation-method}

% TODO: Complete this RES-ML Domain-specific subsection.

\subsection{Evaluation Metrics}
\label{sec:res-ml-rol-evaluation-metrics}

% TODO: Complete this RES-ML Domain-specific subsection.

\subsection{Generalisation and Uncertainty}
\label{sec:res-ml-rol-generalisation-uncertainty}

% TODO: Complete this RES-ML Domain-specific subsection.

\subsection{Simulation and Optimisation Integration}
\label{sec:res-ml-rol-simulation-optimisation-integration}

% TODO: Complete this RES-ML Domain-specific subsection.

\subsection{Candidate Approaches}
\label{sec:res-ml-rol-candidate-approaches}

% TODO: Define the credible candidate approaches considered.

\subsection{Critical Comparison}
\label{sec:res-ml-rol-critical-comparison}

% TODO: Compare candidates using physically and project-relevant criteria.
% Do not select an approach solely on popularity or implementation ease.

\subsection{Assumptions and Applicability}
\label{sec:res-ml-rol-assumptions}

% TODO: State assumptions and the conditions under which they are valid.

\subsection{Limitations and Failure Conditions}
\label{sec:res-ml-rol-limitations}

% TODO: State where the evidence, model or selected approach ceases to be
% reliable or applicable.

\subsection{Project-Specific Conclusions}
\label{sec:res-ml-rol-conclusions}

The active project has no ML training role. A future surrogate may
screen candidate options and identify cases for further simulation,
but it shall not autonomously declare a defence safe or replace
validated hydrodynamic and structural evidence.

\subsection{Selected Approach and Rationale}
\label{sec:res-ml-rol-selected-approach}

% TODO: Record the selected approach only after sufficient evidence exists.
% State the alternatives rejected and the reasons for rejection.

\subsection{Unresolved Decisions}
\label{sec:res-ml-rol-unresolved-decisions}

Unresolved questions are the validated output dataset, accepted
hydrodynamic labels, future structural/damage labels, and the
threshold at which surrogate uncertainty would trigger direct
simulation.

\subsection{Requirements and Handoffs}
\label{sec:res-ml-rol-requirements}

% TODO: State requirements passed to other Research Domains, Software,
% verification activities or later execution work.

\subsection{Deliverable Completion Record}
\label{sec:res-ml-rol-completion-record}

% TODO: Record whether the Deliverable is:
%
% - Not started
% - In progress
% - Draft complete
% - Reviewed
% - Accepted
%
% Acceptance requires:
%
% - the research questions to be answered;
% - claims to be supported by appropriate evidence;
% - assumptions and limitations to be explicit;
% - the project-specific conclusion to be stated;
% - downstream requirements to be traceable;
% - unresolved decisions to be closed or formally deferred.

% -------------------------------------------------------------------------
% WBS ID: RES-ML-DAT
% Deliverable: Learning Dataset and Simulation-Data Requirements
% Status: In progress
% Evidence status: ML dataset deferred by current hydrodynamic scope
% Review status: Not reviewed
% Last updated: 2026-07-19
% Dependencies: RES-DAT-OBS, RES-DAT-SCN, RES-DAT-CAS, RES-GEO-SPEC, RES-VER-SPEC
% Downstream interfaces: RES-ML-REP, RES-ML-TRN, RES-ML-SUR
% -------------------------------------------------------------------------

\section{RES-ML-DAT --- Learning Dataset and Simulation-Data Requirements}
\label{sec:res-ml-dat}

\subsection{Purpose and Scope}
\label{sec:res-ml-dat-purpose}

This Deliverable records future learning-dataset requirements only.
The active project builds hydrodynamic source, corridor, geometry and
output manifests; it does not create an ML training dataset. A future
dataset sample would be a complete simulation case, not an extracted
time step. The deferred dataset relationship is:
\begin{align}
  &\text{corridor identity}
  + \text{fixed-rigid barrier geometry}
  + \text{validated tsunami condition}
  + \text{simulation fidelity and numerical configuration}
  \notag\\
  &\qquad \rightarrow
  \text{hydraulic response}
  + \text{uncertainty and quality state}
  \label{eq:res-ml-dat-sample-definition}
\end{align}

Historical pre-event and current or reconstructed geometries may define
realistic intervention features and contextual or weak labels, but they shall
not be treated as direct causal ground truth for optimal reinforcement.

\subsection{Research Questions}
\label{sec:res-ml-dat-research-questions}

% TODO: State the questions that must be answered before the Deliverable
% can be accepted.

\subsection{Terminology and Definitions}
\label{sec:res-ml-dat-definitions}

% TODO: Define the technical terminology, symbols, quantities and
% classifications used in this Deliverable.

\subsection{Literature Evidence}
\label{sec:res-ml-dat-literature-evidence}

% TODO: Record evidence from peer-reviewed papers, authoritative datasets,
% standards, technical reports and primary documentation.
%
% Do not create a detached literature-summary list. Explain how each source
% supports, contradicts or limits a project decision.

\subsection{Governing Relations and Quantitative Definitions}
\label{sec:res-ml-dat-governing-relations}

% TODO: Record relevant equations, dimensionless groups, empirical models,
% extraction rules or quantitative definitions.
%
% Use align environments for displayed equations.
% Every displayed equation must have a unique \label{eq:...}.
% Do not place punctuation after displayed equations.

\subsection{Machine-Learning Role}
\label{sec:res-ml-dat-machine-learning-role}

% TODO: Complete this RES-ML Domain-specific subsection.

\subsection{Non-ML Baseline}
\label{sec:res-ml-dat-non-ml-baseline}

% TODO: Complete this RES-ML Domain-specific subsection.

\subsection{Input and Output Representation}
\label{sec:res-ml-dat-input-output-representation}

% TODO: Complete this RES-ML Domain-specific subsection.

\subsection{Dataset Requirements}
\label{sec:res-ml-dat-dataset-requirements}

Future ML data groups may include:
\begin{itemize}
  \item corridor, geometry and case identity;
  \item fixed-rigid barrier geometry;
  \item local tsunami-condition descriptors;
  \item fluid-response outputs;
  \item simulation metadata;
  \item convergence state;
  \item physical-validity state;
  \item uncertainty state;
  \item future structural, damage or adequacy labels only if those
    extensions are activated and validated.
\end{itemize}

No ML training dataset may be accepted before the selected event,
local hydrodynamic model and impact/loading outputs have passed the
V1--V2--V3 validation gates.

\subsection{Model or Architecture Candidates}
\label{sec:res-ml-dat-model-architecture-candidates}

% TODO: Complete this RES-ML Domain-specific subsection.

\subsection{Training and Validation Method}
\label{sec:res-ml-dat-training-validation-method}

The previous mandatory secondary-event holdout requirement remains
removed from the active project. Future ML validation would require
grouped partitioning based on complete simulation cases, held-out
geometry configurations, held-out condition ranges and independent
final test simulations.

Time steps from the same simulation shall never be randomly split between
future training and test datasets.

\subsection{Evaluation Metrics}
\label{sec:res-ml-dat-evaluation-metrics}

% TODO: Complete this RES-ML Domain-specific subsection.

\subsection{Generalisation and Uncertainty}
\label{sec:res-ml-dat-generalisation-uncertainty}

% TODO: Complete this RES-ML Domain-specific subsection.

\subsection{Simulation and Optimisation Integration}
\label{sec:res-ml-dat-simulation-optimisation-integration}

% TODO: Complete this RES-ML Domain-specific subsection.

\subsection{Candidate Approaches}
\label{sec:res-ml-dat-candidate-approaches}

% TODO: Define the credible candidate approaches considered.

\subsection{Critical Comparison}
\label{sec:res-ml-dat-critical-comparison}

% TODO: Compare candidates using physically and project-relevant criteria.
% Do not select an approach solely on popularity or implementation ease.

\subsection{Assumptions and Applicability}
\label{sec:res-ml-dat-assumptions}

% TODO: State assumptions and the conditions under which they are valid.

\subsection{Limitations and Failure Conditions}
\label{sec:res-ml-dat-limitations}

% TODO: State where the evidence, model or selected approach ceases to be
% reliable or applicable.

\subsection{Project-Specific Conclusions}
\label{sec:res-ml-dat-conclusions}

The initial learning dataset is a controlled simulation-derived dataset for one
validated event-conditioned domain. It is not a broad multi-event dataset and
does not support cross-event generalisation claims.

\subsection{Selected Approach and Rationale}
\label{sec:res-ml-dat-selected-approach}

% TODO: Record the selected approach only after sufficient evidence exists.
% State the alternatives rejected and the reasons for rejection.

\subsection{Unresolved Decisions}
\label{sec:res-ml-dat-unresolved-decisions}

% TODO: Record unresolved questions, required evidence and decision owners.

\subsection{Requirements and Handoffs}
\label{sec:res-ml-dat-requirements}

% TODO: State requirements passed to other Research Domains, Software,
% verification activities or later execution work.

\subsection{Deliverable Completion Record}
\label{sec:res-ml-dat-completion-record}

% TODO: Record whether the Deliverable is:
%
% - Not started
% - In progress
% - Draft complete
% - Reviewed
% - Accepted
%
% Acceptance requires:
%
% - the research questions to be answered;
% - claims to be supported by appropriate evidence;
% - assumptions and limitations to be explicit;
% - the project-specific conclusion to be stated;
% - downstream requirements to be traceable;
% - unresolved decisions to be closed or formally deferred.

% -------------------------------------------------------------------------
% WBS ID: RES-ML-REP
% Deliverable: Input, Output and Geometry Representation
% Status: In progress
% Evidence status: ML representation deferred; geometry manifest remains active
% Review status: Not reviewed
% Last updated: 2026-07-19
% Dependencies: RES-GEO-PAR, RES-GEO-MET, RES-GEO-REP, RES-DAT-SCN
% Downstream interfaces: RES-ML-DAT, RES-ML-ALG, RES-ML-SUR
% -------------------------------------------------------------------------

\section{RES-ML-REP --- Input, Output and Geometry Representation}
\label{sec:res-ml-rep}

\subsection{Purpose and Scope}
\label{sec:res-ml-rep-purpose}

This Deliverable records future ML representation requirements. The
active representation work is the RES-GEO/RES-DAT manifest for fixed-
rigid barrier-class hydrodynamic comparison. Raw voxel, mesh, graph,
full-field learning and learned geometry representations are deferred.

\subsection{Research Questions}
\label{sec:res-ml-rep-research-questions}

% TODO: State the questions that must be answered before the Deliverable
% can be accepted.

\subsection{Terminology and Definitions}
\label{sec:res-ml-rep-definitions}

% TODO: Define the technical terminology, symbols, quantities and
% classifications used in this Deliverable.

\subsection{Literature Evidence}
\label{sec:res-ml-rep-literature-evidence}

% TODO: Record evidence from peer-reviewed papers, authoritative datasets,
% standards, technical reports and primary documentation.
%
% Do not create a detached literature-summary list. Explain how each source
% supports, contradicts or limits a project decision.

\subsection{Governing Relations and Quantitative Definitions}
\label{sec:res-ml-rep-governing-relations}

% TODO: Record relevant equations, dimensionless groups, empirical models,
% extraction rules or quantitative definitions.
%
% Use align environments for displayed equations.
% Every displayed equation must have a unique \label{eq:...}.
% Do not place punctuation after displayed equations.

\subsection{Machine-Learning Role}
\label{sec:res-ml-rep-machine-learning-role}

% TODO: Complete this RES-ML Domain-specific subsection.

\subsection{Non-ML Baseline}
\label{sec:res-ml-rep-non-ml-baseline}

% TODO: Complete this RES-ML Domain-specific subsection.

\subsection{Input and Output Representation}
\label{sec:res-ml-rep-input-output-representation}

Future input descriptors may include corridor identity, barrier class,
height, width and thickness, orientation, face angle, resolved openings,
incident water depth, incident velocity, wave direction, duration or
time-history descriptors. Material, foundation and reinforcement
descriptors are excluded until structural or optimisation extensions
are activated.

Future outputs may include peak force, peak overturning moment,
overtopping, run-up and energy-transfer or dissipation metrics.
Displacement, stress, strain, permanent deformation, adequacy class,
damage class and structural margin are deferred.

\subsection{Dataset Requirements}
\label{sec:res-ml-rep-dataset-requirements}

% TODO: Complete this RES-ML Domain-specific subsection.

\subsection{Model or Architecture Candidates}
\label{sec:res-ml-rep-model-architecture-candidates}

% TODO: Complete this RES-ML Domain-specific subsection.

\subsection{Training and Validation Method}
\label{sec:res-ml-rep-training-validation-method}

% TODO: Complete this RES-ML Domain-specific subsection.

\subsection{Evaluation Metrics}
\label{sec:res-ml-rep-evaluation-metrics}

% TODO: Complete this RES-ML Domain-specific subsection.

\subsection{Generalisation and Uncertainty}
\label{sec:res-ml-rep-generalisation-uncertainty}

Every representation shall carry feature provenance and known uncertainty so
that missing geometry, uncertain material properties and inferred foundation
conditions are not silently treated as measured quantities.

\subsection{Simulation and Optimisation Integration}
\label{sec:res-ml-rep-simulation-optimisation-integration}

% TODO: Complete this RES-ML Domain-specific subsection.

\subsection{Candidate Approaches}
\label{sec:res-ml-rep-candidate-approaches}

% TODO: Define the credible candidate approaches considered.

\subsection{Critical Comparison}
\label{sec:res-ml-rep-critical-comparison}

% TODO: Compare candidates using physically and project-relevant criteria.
% Do not select an approach solely on popularity or implementation ease.

\subsection{Assumptions and Applicability}
\label{sec:res-ml-rep-assumptions}

% TODO: State assumptions and the conditions under which they are valid.

\subsection{Limitations and Failure Conditions}
\label{sec:res-ml-rep-limitations}

% TODO: State where the evidence, model or selected approach ceases to be
% reliable or applicable.

\subsection{Project-Specific Conclusions}
\label{sec:res-ml-rep-conclusions}

The current representation family is a physics-informed manifest for
fixed-rigid hydrodynamic comparison. It may later become a tabular ML
descriptor set, but it is not currently a surrogate input, generative
geometry encoding or unrestricted design-space representation.

\subsection{Selected Approach and Rationale}
\label{sec:res-ml-rep-selected-approach}

% TODO: Record the selected approach only after sufficient evidence exists.
% State the alternatives rejected and the reasons for rejection.

\subsection{Unresolved Decisions}
\label{sec:res-ml-rep-unresolved-decisions}

% TODO: Record unresolved questions, required evidence and decision owners.

\subsection{Requirements and Handoffs}
\label{sec:res-ml-rep-requirements}

% TODO: State requirements passed to other Research Domains, Software,
% verification activities or later execution work.

\subsection{Deliverable Completion Record}
\label{sec:res-ml-rep-completion-record}

% TODO: Record whether the Deliverable is:
%
% - Not started
% - In progress
% - Draft complete
% - Reviewed
% - Accepted
%
% Acceptance requires:
%
% - the research questions to be answered;
% - claims to be supported by appropriate evidence;
% - assumptions and limitations to be explicit;
% - the project-specific conclusion to be stated;
% - downstream requirements to be traceable;
% - unresolved decisions to be closed or formally deferred.

% -------------------------------------------------------------------------
% WBS ID: RES-ML-ALG
% Deliverable: Learning Algorithm and Architecture Assessment
% Status: In progress
% Evidence status: Algorithm assessment deferred by current hydrodynamic scope
% Review status: Not reviewed
% Last updated: 2026-07-19
% Dependencies: RES-ML-DAT, RES-ML-REP, RES-ML-TRN
% Downstream interfaces: RES-ML-SUR, RES-ML-SPEC
% -------------------------------------------------------------------------

\section{RES-ML-ALG --- Learning Algorithm and Architecture Assessment}
\label{sec:res-ml-alg}

\subsection{Purpose and Scope}
\label{sec:res-ml-alg-purpose}

This Deliverable records future algorithm candidates only. The active
project does not assess or select an ML algorithm; it first produces
validated hydrodynamic outputs and structured provenance.

\subsection{Research Questions}
\label{sec:res-ml-alg-research-questions}

% TODO: State the questions that must be answered before the Deliverable
% can be accepted.

\subsection{Terminology and Definitions}
\label{sec:res-ml-alg-definitions}

% TODO: Define the technical terminology, symbols, quantities and
% classifications used in this Deliverable.

\subsection{Literature Evidence}
\label{sec:res-ml-alg-literature-evidence}

% TODO: Record evidence from peer-reviewed papers, authoritative datasets,
% standards, technical reports and primary documentation.
%
% Do not create a detached literature-summary list. Explain how each source
% supports, contradicts or limits a project decision.

The Intelligent Systems lecture material may be used as terminology and
foundation context for linear regression, polynomial regression, logistic
regression, Naive Bayes, model selection, kernels, support-vector methods,
clustering, dimensionality reduction, neural networks, linear algebra, calculus,
probability and statistics. It shall not be used as the sole evidence for
research claims that require peer-reviewed or authoritative primary sources.

\subsection{Governing Relations and Quantitative Definitions}
\label{sec:res-ml-alg-governing-relations}

% TODO: Record relevant equations, dimensionless groups, empirical models,
% extraction rules or quantitative definitions.
%
% Use align environments for displayed equations.
% Every displayed equation must have a unique \label{eq:...}.
% Do not place punctuation after displayed equations.

\subsection{Machine-Learning Role}
\label{sec:res-ml-alg-machine-learning-role}

% TODO: Complete this RES-ML Domain-specific subsection.

\subsection{Non-ML Baseline}
\label{sec:res-ml-alg-non-ml-baseline}

% TODO: Complete this RES-ML Domain-specific subsection.

\subsection{Input and Output Representation}
\label{sec:res-ml-alg-input-output-representation}

% TODO: Complete this RES-ML Domain-specific subsection.

\subsection{Dataset Requirements}
\label{sec:res-ml-alg-dataset-requirements}

% TODO: Complete this RES-ML Domain-specific subsection.

\subsection{Model or Architecture Candidates}
\label{sec:res-ml-alg-model-architecture-candidates}

If ML is reactivated later, the assessment hierarchy may be:
\begin{enumerate}
  \item linear regression baseline;
  \item regularised polynomial regression;
  \item logistic regression for classification;
  \item linear and nonlinear support-vector methods;
  \item kernel-based regression;
  \item tree ensembles or boosting;
  \item Gaussian-process methods where computationally feasible;
  \item feed-forward neural networks only when dataset volume supports them;
  \item PCA or another controlled dimensionality-reduction method for
    correlated outputs or histories;
  \item clustering for exploratory analysis only.
\end{enumerate}

\subsection{Training and Validation Method}
\label{sec:res-ml-alg-training-validation-method}

% TODO: Complete this RES-ML Domain-specific subsection.

\subsection{Evaluation Metrics}
\label{sec:res-ml-alg-evaluation-metrics}

% TODO: Complete this RES-ML Domain-specific subsection.

\subsection{Generalisation and Uncertainty}
\label{sec:res-ml-alg-generalisation-uncertainty}

% TODO: Complete this RES-ML Domain-specific subsection.

\subsection{Simulation and Optimisation Integration}
\label{sec:res-ml-alg-simulation-optimisation-integration}

% TODO: Complete this RES-ML Domain-specific subsection.

\subsection{Candidate Approaches}
\label{sec:res-ml-alg-candidate-approaches}

% TODO: Define the credible candidate approaches considered.

\subsection{Critical Comparison}
\label{sec:res-ml-alg-critical-comparison}

Model selection shall consider dataset size, dimensionality, nonlinearity,
uncertainty, interpretability, extrapolation, failure-boundary behaviour and
implementation cost. Popularity or implementation convenience alone is
insufficient evidence for selection.

\subsection{Assumptions and Applicability}
\label{sec:res-ml-alg-assumptions}

% TODO: State assumptions and the conditions under which they are valid.

\subsection{Limitations and Failure Conditions}
\label{sec:res-ml-alg-limitations}

% TODO: State where the evidence, model or selected approach ceases to be
% reliable or applicable.

\subsection{Project-Specific Conclusions}
\label{sec:res-ml-alg-conclusions}

The initial algorithm strategy is baseline-first and uncertainty-aware. A final
algorithm remains unresolved until data volume, feature dimensionality and
validation behaviour are known.

\subsection{Selected Approach and Rationale}
\label{sec:res-ml-alg-selected-approach}

% TODO: Record the selected approach only after sufficient evidence exists.
% State the alternatives rejected and the reasons for rejection.

\subsection{Unresolved Decisions}
\label{sec:res-ml-alg-unresolved-decisions}

% TODO: Record unresolved questions, required evidence and decision owners.

\subsection{Requirements and Handoffs}
\label{sec:res-ml-alg-requirements}

% TODO: State requirements passed to other Research Domains, Software,
% verification activities or later execution work.

\subsection{Deliverable Completion Record}
\label{sec:res-ml-alg-completion-record}

% TODO: Record whether the Deliverable is:
%
% - Not started
% - In progress
% - Draft complete
% - Reviewed
% - Accepted
%
% Acceptance requires:
%
% - the research questions to be answered;
% - claims to be supported by appropriate evidence;
% - assumptions and limitations to be explicit;
% - the project-specific conclusion to be stated;
% - downstream requirements to be traceable;
% - unresolved decisions to be closed or formally deferred.

% -------------------------------------------------------------------------
% WBS ID: RES-ML-OBJ
% Deliverable: Loss Functions, Rewards and Physical Constraints
% Status: In progress
% Evidence status: Not assessed
% Review status: Not reviewed
% Last updated: 2026-06-30
% Dependencies: RES-ML-DAT, RES-ML-ALG, RES-VER-MET
% Downstream interfaces: RES-ML-TRN, RES-ML-SUR, RES-ML-SPEC
% -------------------------------------------------------------------------

\section{RES-ML-OBJ --- Loss Functions, Rewards and Physical Constraints}
\label{sec:res-ml-obj}

\subsection{Purpose and Scope}
\label{sec:res-ml-obj-purpose}

This Deliverable defines prediction losses, classification labels and
engineering constraints for the event-conditioned supervised-learning task.
Reinforcement-learning rewards are removed from the active scope.

\subsection{Research Questions}
\label{sec:res-ml-obj-research-questions}

% TODO: State the questions that must be answered before the Deliverable
% can be accepted.

\subsection{Terminology and Definitions}
\label{sec:res-ml-obj-definitions}

% TODO: Define the technical terminology, symbols, quantities and
% classifications used in this Deliverable.

\subsection{Literature Evidence}
\label{sec:res-ml-obj-literature-evidence}

% TODO: Record evidence from peer-reviewed papers, authoritative datasets,
% standards, technical reports and primary documentation.
%
% Do not create a detached literature-summary list. Explain how each source
% supports, contradicts or limits a project decision.

\subsection{Governing Relations and Quantitative Definitions}
\label{sec:res-ml-obj-governing-relations}

% TODO: Record relevant equations, dimensionless groups, empirical models,
% extraction rules or quantitative definitions.
%
% Use align environments for displayed equations.
% Every displayed equation must have a unique \label{eq:...}.
% Do not place punctuation after displayed equations.

\subsection{Machine-Learning Role}
\label{sec:res-ml-obj-machine-learning-role}

% TODO: Complete this RES-ML Domain-specific subsection.

\subsection{Non-ML Baseline}
\label{sec:res-ml-obj-non-ml-baseline}

% TODO: Complete this RES-ML Domain-specific subsection.

\subsection{Input and Output Representation}
\label{sec:res-ml-obj-input-output-representation}

% TODO: Complete this RES-ML Domain-specific subsection.

\subsection{Dataset Requirements}
\label{sec:res-ml-obj-dataset-requirements}

% TODO: Complete this RES-ML Domain-specific subsection.

\subsection{Model or Architecture Candidates}
\label{sec:res-ml-obj-model-architecture-candidates}

% TODO: Complete this RES-ML Domain-specific subsection.

\subsection{Training and Validation Method}
\label{sec:res-ml-obj-training-validation-method}

% TODO: Complete this RES-ML Domain-specific subsection.

\subsection{Evaluation Metrics}
\label{sec:res-ml-obj-evaluation-metrics}

The evaluation methodology shall include regression losses, classification
losses, multi-output normalisation, regularisation, class imbalance handling,
physical bounds, invalid prediction handling, and asymmetric consequences of
classification errors.

The following error asymmetry is mandatory:
\begin{itemize}
  \item unsafe defence predicted safe: high-consequence error;
  \item safe defence predicted unsafe: conservative but inefficient error.
\end{itemize}

Acceptance criteria and model-selection metrics shall account for this
asymmetry rather than relying only on aggregate accuracy.

\subsection{Generalisation and Uncertainty}
\label{sec:res-ml-obj-generalisation-uncertainty}

% TODO: Complete this RES-ML Domain-specific subsection.

\subsection{Simulation and Optimisation Integration}
\label{sec:res-ml-obj-simulation-optimisation-integration}

Losses and labels shall support reinforcement screening only within the
validated event-conditioned domain. Queries outside the domain or with invalid
physical predictions shall be rejected or routed to direct simulation.

\subsection{Candidate Approaches}
\label{sec:res-ml-obj-candidate-approaches}

% TODO: Define the credible candidate approaches considered.

\subsection{Critical Comparison}
\label{sec:res-ml-obj-critical-comparison}

% TODO: Compare candidates using physically and project-relevant criteria.
% Do not select an approach solely on popularity or implementation ease.

\subsection{Assumptions and Applicability}
\label{sec:res-ml-obj-assumptions}

% TODO: State assumptions and the conditions under which they are valid.

\subsection{Limitations and Failure Conditions}
\label{sec:res-ml-obj-limitations}

% TODO: State where the evidence, model or selected approach ceases to be
% reliable or applicable.

\subsection{Project-Specific Conclusions}
\label{sec:res-ml-obj-conclusions}

The project uses supervised response and adequacy targets, not
reinforcement-learning reward optimisation. Safety-relevant false-safe
predictions dominate the risk model.

\subsection{Selected Approach and Rationale}
\label{sec:res-ml-obj-selected-approach}

% TODO: Record the selected approach only after sufficient evidence exists.
% State the alternatives rejected and the reasons for rejection.

\subsection{Unresolved Decisions}
\label{sec:res-ml-obj-unresolved-decisions}

% TODO: Record unresolved questions, required evidence and decision owners.

\subsection{Requirements and Handoffs}
\label{sec:res-ml-obj-requirements}

% TODO: State requirements passed to other Research Domains, Software,
% verification activities or later execution work.

\subsection{Deliverable Completion Record}
\label{sec:res-ml-obj-completion-record}

% TODO: Record whether the Deliverable is:
%
% - Not started
% - In progress
% - Draft complete
% - Reviewed
% - Accepted
%
% Acceptance requires:
%
% - the research questions to be answered;
% - claims to be supported by appropriate evidence;
% - assumptions and limitations to be explicit;
% - the project-specific conclusion to be stated;
% - downstream requirements to be traceable;
% - unresolved decisions to be closed or formally deferred.

% -------------------------------------------------------------------------
% WBS ID: RES-ML-SUR
% Deliverable: Surrogate, Multifidelity and Active-Learning Strategy
% Status: In progress
% Evidence status: Surrogate and active learning deferred
% Review status: Not reviewed
% Last updated: 2026-07-19
% Dependencies: RES-ML-DAT, RES-ML-ALG, RES-ML-TRN, RES-VER-SPEC
% Downstream interfaces: RES-ML-INT, RES-OPT-DOE, RES-OPT-SPEC
% -------------------------------------------------------------------------

\section{RES-ML-SUR --- Surrogate, Multifidelity and Active-Learning Strategy}
\label{sec:res-ml-sur}

\subsection{Purpose and Scope}
\label{sec:res-ml-sur-purpose}

This Deliverable records surrogate, multifidelity and active-learning
ideas as deferred future work. The active project does not train a
surrogate or run active sampling; it prepares validated hydrodynamic
outputs for possible future use.

\subsection{Research Questions}
\label{sec:res-ml-sur-research-questions}

% TODO: State the questions that must be answered before the Deliverable
% can be accepted.

\subsection{Terminology and Definitions}
\label{sec:res-ml-sur-definitions}

% TODO: Define the technical terminology, symbols, quantities and
% classifications used in this Deliverable.

\subsection{Literature Evidence}
\label{sec:res-ml-sur-literature-evidence}

% TODO: Record evidence from peer-reviewed papers, authoritative datasets,
% standards, technical reports and primary documentation.
%
% Do not create a detached literature-summary list. Explain how each source
% supports, contradicts or limits a project decision.

\subsection{Governing Relations and Quantitative Definitions}
\label{sec:res-ml-sur-governing-relations}

% TODO: Record relevant equations, dimensionless groups, empirical models,
% extraction rules or quantitative definitions.
%
% Use align environments for displayed equations.
% Every displayed equation must have a unique \label{eq:...}.
% Do not place punctuation after displayed equations.

\subsection{Machine-Learning Role}
\label{sec:res-ml-sur-machine-learning-role}

The deferred future query is:
\begin{align}
  &\text{corridor}
  + \text{validated-event condition}
  + \text{fixed-rigid barrier class}
  \notag\\
  &\qquad \rightarrow
  \text{predicted response}
  + \text{predictive uncertainty}
  \label{eq:res-ml-sur-query}
\end{align}

Different tsunami conditions initially mean controlled variations within a
defensible envelope around the validated event, such as local wave-height
variation, inundation-depth variation, velocity variation, direction variation
and selected severity scaling. This shall not be described as generalisation to
unrelated earthquakes, coastlines or tsunami events.

\subsection{Non-ML Baseline}
\label{sec:res-ml-sur-non-ml-baseline}

% TODO: Complete this RES-ML Domain-specific subsection.

\subsection{Input and Output Representation}
\label{sec:res-ml-sur-input-output-representation}

% TODO: Complete this RES-ML Domain-specific subsection.

\subsection{Dataset Requirements}
\label{sec:res-ml-sur-dataset-requirements}

% TODO: Complete this RES-ML Domain-specific subsection.

\subsection{Model or Architecture Candidates}
\label{sec:res-ml-sur-model-architecture-candidates}

% TODO: Complete this RES-ML Domain-specific subsection.

\subsection{Training and Validation Method}
\label{sec:res-ml-sur-training-validation-method}

% TODO: Complete this RES-ML Domain-specific subsection.

\subsection{Evaluation Metrics}
\label{sec:res-ml-sur-evaluation-metrics}

% TODO: Complete this RES-ML Domain-specific subsection.

\subsection{Generalisation and Uncertainty}
\label{sec:res-ml-sur-generalisation-uncertainty}

% TODO: Complete this RES-ML Domain-specific subsection.

\subsection{Simulation and Optimisation Integration}
\label{sec:res-ml-sur-simulation-optimisation-integration}

A future active-sampling strategy may prioritise high-uncertainty
hydrodynamic regions, underrepresented geometry configurations and
cases where model disagreement is high. Failure-boundary and direct
FSI confirmation cases remain deferred until structural/FSI extensions
exist.

\subsection{Candidate Approaches}
\label{sec:res-ml-sur-candidate-approaches}

% TODO: Define the credible candidate approaches considered.

\subsection{Critical Comparison}
\label{sec:res-ml-sur-critical-comparison}

% TODO: Compare candidates using physically and project-relevant criteria.
% Do not select an approach solely on popularity or implementation ease.

\subsection{Assumptions and Applicability}
\label{sec:res-ml-sur-assumptions}

% TODO: State assumptions and the conditions under which they are valid.

\subsection{Limitations and Failure Conditions}
\label{sec:res-ml-sur-limitations}

% TODO: State where the evidence, model or selected approach ceases to be
% reliable or applicable.

\subsection{Project-Specific Conclusions}
\label{sec:res-ml-sur-conclusions}

A future surrogate may be an event-conditioned screening tool. Direct
simulation remains the source of accepted physical response and is
required whenever uncertainty, domain limits or consequence level
demand confirmation.

\subsection{Selected Approach and Rationale}
\label{sec:res-ml-sur-selected-approach}

% TODO: Record the selected approach only after sufficient evidence exists.
% State the alternatives rejected and the reasons for rejection.

\subsection{Unresolved Decisions}
\label{sec:res-ml-sur-unresolved-decisions}

% TODO: Record unresolved questions, required evidence and decision owners.

\subsection{Requirements and Handoffs}
\label{sec:res-ml-sur-requirements}

% TODO: State requirements passed to other Research Domains, Software,
% verification activities or later execution work.

\subsection{Deliverable Completion Record}
\label{sec:res-ml-sur-completion-record}

% TODO: Record whether the Deliverable is:
%
% - Not started
% - In progress
% - Draft complete
% - Reviewed
% - Accepted
%
% Acceptance requires:
%
% - the research questions to be answered;
% - claims to be supported by appropriate evidence;
% - assumptions and limitations to be explicit;
% - the project-specific conclusion to be stated;
% - downstream requirements to be traceable;
% - unresolved decisions to be closed or formally deferred.

% -------------------------------------------------------------------------
% WBS ID: RES-ML-TRN
% Deliverable: Training, Validation and Experiment Methodology
% Status: In progress
% Evidence status: ML training deferred by current hydrodynamic scope
% Review status: Not reviewed
% Last updated: 2026-07-19
% Dependencies: RES-ML-DAT, RES-ML-ALG, RES-ML-OBJ
% Downstream interfaces: RES-ML-GEN, RES-ML-SPEC
% -------------------------------------------------------------------------

\section{RES-ML-TRN --- Training, Validation and Experiment Methodology}
\label{sec:res-ml-trn}

\subsection{Purpose and Scope}
\label{sec:res-ml-trn-purpose}

This Deliverable records future training, grouped validation and
model-selection rules only. There is no active ML training in the
current hydrodynamic Research phase. If ML is reactivated later, the
simulation run shall be the indivisible grouping unit.

\subsection{Research Questions}
\label{sec:res-ml-trn-research-questions}

% TODO: State the questions that must be answered before the Deliverable
% can be accepted.

\subsection{Terminology and Definitions}
\label{sec:res-ml-trn-definitions}

% TODO: Define the technical terminology, symbols, quantities and
% classifications used in this Deliverable.

\subsection{Literature Evidence}
\label{sec:res-ml-trn-literature-evidence}

% TODO: Record evidence from peer-reviewed papers, authoritative datasets,
% standards, technical reports and primary documentation.
%
% Do not create a detached literature-summary list. Explain how each source
% supports, contradicts or limits a project decision.

\subsection{Governing Relations and Quantitative Definitions}
\label{sec:res-ml-trn-governing-relations}

% TODO: Record relevant equations, dimensionless groups, empirical models,
% extraction rules or quantitative definitions.
%
% Use align environments for displayed equations.
% Every displayed equation must have a unique \label{eq:...}.
% Do not place punctuation after displayed equations.

\subsection{Machine-Learning Role}
\label{sec:res-ml-trn-machine-learning-role}

% TODO: Complete this RES-ML Domain-specific subsection.

\subsection{Non-ML Baseline}
\label{sec:res-ml-trn-non-ml-baseline}

% TODO: Complete this RES-ML Domain-specific subsection.

\subsection{Input and Output Representation}
\label{sec:res-ml-trn-input-output-representation}

% TODO: Complete this RES-ML Domain-specific subsection.

\subsection{Dataset Requirements}
\label{sec:res-ml-trn-dataset-requirements}

% TODO: Complete this RES-ML Domain-specific subsection.

\subsection{Model or Architecture Candidates}
\label{sec:res-ml-trn-model-architecture-candidates}

% TODO: Complete this RES-ML Domain-specific subsection.

\subsection{Training and Validation Method}
\label{sec:res-ml-trn-training-validation-method}

If ML is reactivated later, the previous secondary-event holdout
requirement shall be replaced with:
\begin{itemize}
  \item grouped splits by complete simulation;
  \item held-out geometry or intervention options;
  \item held-out condition ranges;
  \item structure-level holdout where sufficient structures exist;
  \item grouped or nested cross-validation;
  \item fixed seeds;
  \item non-ML baselines;
  \item independent final test cases.
\end{itemize}

Time steps or spatial samples from one simulation shall not be randomly divided
between future training and test partitions.

\subsection{Evaluation Metrics}
\label{sec:res-ml-trn-evaluation-metrics}

% TODO: Complete this RES-ML Domain-specific subsection.

\subsection{Generalisation and Uncertainty}
\label{sec:res-ml-trn-generalisation-uncertainty}

% TODO: Complete this RES-ML Domain-specific subsection.

\subsection{Simulation and Optimisation Integration}
\label{sec:res-ml-trn-simulation-optimisation-integration}

% TODO: Complete this RES-ML Domain-specific subsection.

\subsection{Candidate Approaches}
\label{sec:res-ml-trn-candidate-approaches}

% TODO: Define the credible candidate approaches considered.

\subsection{Critical Comparison}
\label{sec:res-ml-trn-critical-comparison}

% TODO: Compare candidates using physically and project-relevant criteria.
% Do not select an approach solely on popularity or implementation ease.

\subsection{Assumptions and Applicability}
\label{sec:res-ml-trn-assumptions}

% TODO: State assumptions and the conditions under which they are valid.

\subsection{Limitations and Failure Conditions}
\label{sec:res-ml-trn-limitations}

% TODO: State where the evidence, model or selected approach ceases to be
% reliable or applicable.

\subsection{Project-Specific Conclusions}
\label{sec:res-ml-trn-conclusions}

Future ML validation shall be grouped by engineering case and
geometry/intervention option, not by arbitrary sample rows. This
protects against leakage from repeated outputs of the same simulation
campaign.

\subsection{Selected Approach and Rationale}
\label{sec:res-ml-trn-selected-approach}

% TODO: Record the selected approach only after sufficient evidence exists.
% State the alternatives rejected and the reasons for rejection.

\subsection{Unresolved Decisions}
\label{sec:res-ml-trn-unresolved-decisions}

% TODO: Record unresolved questions, required evidence and decision owners.

\subsection{Requirements and Handoffs}
\label{sec:res-ml-trn-requirements}

% TODO: State requirements passed to other Research Domains, Software,
% verification activities or later execution work.

\subsection{Deliverable Completion Record}
\label{sec:res-ml-trn-completion-record}

% TODO: Record whether the Deliverable is:
%
% - Not started
% - In progress
% - Draft complete
% - Reviewed
% - Accepted
%
% Acceptance requires:
%
% - the research questions to be answered;
% - claims to be supported by appropriate evidence;
% - assumptions and limitations to be explicit;
% - the project-specific conclusion to be stated;
% - downstream requirements to be traceable;
% - unresolved decisions to be closed or formally deferred.

% -------------------------------------------------------------------------
% WBS ID: RES-ML-GEN
% Deliverable: Generalisation, Uncertainty and Interpretability
% Status: In progress
% Evidence status: Not assessed
% Review status: Not reviewed
% Last updated: 2026-06-30
% Dependencies: RES-ML-TRN, RES-ML-SUR, RES-DAT-CAS
% Downstream interfaces: RES-ML-INT, RES-ML-SPEC
% -------------------------------------------------------------------------

\section{RES-ML-GEN --- Generalisation, Uncertainty and Interpretability}
\label{sec:res-ml-gen}

\subsection{Purpose and Scope}
\label{sec:res-ml-gen-purpose}

This Deliverable restricts generalisation claims to the defined
event-conditioned domain. Cross-event performance is removed as an initial
acceptance requirement.

\subsection{Research Questions}
\label{sec:res-ml-gen-research-questions}

% TODO: State the questions that must be answered before the Deliverable
% can be accepted.

\subsection{Terminology and Definitions}
\label{sec:res-ml-gen-definitions}

% TODO: Define the technical terminology, symbols, quantities and
% classifications used in this Deliverable.

\subsection{Literature Evidence}
\label{sec:res-ml-gen-literature-evidence}

% TODO: Record evidence from peer-reviewed papers, authoritative datasets,
% standards, technical reports and primary documentation.
%
% Do not create a detached literature-summary list. Explain how each source
% supports, contradicts or limits a project decision.

\subsection{Governing Relations and Quantitative Definitions}
\label{sec:res-ml-gen-governing-relations}

% TODO: Record relevant equations, dimensionless groups, empirical models,
% extraction rules or quantitative definitions.
%
% Use align environments for displayed equations.
% Every displayed equation must have a unique \label{eq:...}.
% Do not place punctuation after displayed equations.

\subsection{Machine-Learning Role}
\label{sec:res-ml-gen-machine-learning-role}

% TODO: Complete this RES-ML Domain-specific subsection.

\subsection{Non-ML Baseline}
\label{sec:res-ml-gen-non-ml-baseline}

% TODO: Complete this RES-ML Domain-specific subsection.

\subsection{Input and Output Representation}
\label{sec:res-ml-gen-input-output-representation}

% TODO: Complete this RES-ML Domain-specific subsection.

\subsection{Dataset Requirements}
\label{sec:res-ml-gen-dataset-requirements}

% TODO: Complete this RES-ML Domain-specific subsection.

\subsection{Model or Architecture Candidates}
\label{sec:res-ml-gen-model-architecture-candidates}

% TODO: Complete this RES-ML Domain-specific subsection.

\subsection{Training and Validation Method}
\label{sec:res-ml-gen-training-validation-method}

% TODO: Complete this RES-ML Domain-specific subsection.

\subsection{Evaluation Metrics}
\label{sec:res-ml-gen-evaluation-metrics}

% TODO: Complete this RES-ML Domain-specific subsection.

\subsection{Generalisation and Uncertainty}
\label{sec:res-ml-gen-generalisation-uncertainty}

The methodology shall define interpolation limits, extrapolation detection,
held-out intervention tests, held-out severity tests, uncertainty calibration,
out-of-distribution rejection, feature sensitivity, engineering
interpretability, permitted-use boundaries, and conditions requiring a new FSI
simulation.

\subsection{Simulation and Optimisation Integration}
\label{sec:res-ml-gen-simulation-optimisation-integration}

% TODO: Complete this RES-ML Domain-specific subsection.

\subsection{Candidate Approaches}
\label{sec:res-ml-gen-candidate-approaches}

% TODO: Define the credible candidate approaches considered.

\subsection{Critical Comparison}
\label{sec:res-ml-gen-critical-comparison}

% TODO: Compare candidates using physically and project-relevant criteria.
% Do not select an approach solely on popularity or implementation ease.

\subsection{Assumptions and Applicability}
\label{sec:res-ml-gen-assumptions}

% TODO: State assumptions and the conditions under which they are valid.

\subsection{Limitations and Failure Conditions}
\label{sec:res-ml-gen-limitations}

The surrogate ceases to be applicable outside the validated event envelope,
outside represented defence and reinforcement types, where feature provenance is
insufficient, where predictive uncertainty is excessive, or where the
engineering consequence of an incorrect prediction requires direct simulation.

\subsection{Project-Specific Conclusions}
\label{sec:res-ml-gen-conclusions}

The initial project may claim only within-domain event-conditioned performance.
Cross-event generalisation remains future work.

\subsection{Selected Approach and Rationale}
\label{sec:res-ml-gen-selected-approach}

% TODO: Record the selected approach only after sufficient evidence exists.
% State the alternatives rejected and the reasons for rejection.

\subsection{Unresolved Decisions}
\label{sec:res-ml-gen-unresolved-decisions}

% TODO: Record unresolved questions, required evidence and decision owners.

\subsection{Requirements and Handoffs}
\label{sec:res-ml-gen-requirements}

% TODO: State requirements passed to other Research Domains, Software,
% verification activities or later execution work.

\subsection{Deliverable Completion Record}
\label{sec:res-ml-gen-completion-record}

% TODO: Record whether the Deliverable is:
%
% - Not started
% - In progress
% - Draft complete
% - Reviewed
% - Accepted
%
% Acceptance requires:
%
% - the research questions to be answered;
% - claims to be supported by appropriate evidence;
% - assumptions and limitations to be explicit;
% - the project-specific conclusion to be stated;
% - downstream requirements to be traceable;
% - unresolved decisions to be closed or formally deferred.

% -------------------------------------------------------------------------
% WBS ID: RES-ML-INT
% Deliverable: Integration with Simulation and Optimisation
% Status: In progress
% Evidence status: Integration deferred by current hydrodynamic scope
% Review status: Not reviewed
% Last updated: 2026-07-19
% Dependencies: RES-ML-SUR, RES-VER-SPEC, RES-OPT-SPEC
% Downstream interfaces: Software implementation planning, RES-ML-SPEC
% -------------------------------------------------------------------------

\section{RES-ML-INT --- Integration with Simulation and Optimisation}
\label{sec:res-ml-int}

\subsection{Purpose and Scope}
\label{sec:res-ml-int-purpose}

This Deliverable records a deferred integration path. During the
active project phase, Software integration is limited to data
handling, structured outputs and validation reports for the regional
and local hydrodynamic models. Surrogate, optimisation and FSI
integration remain future work.

\subsection{Research Questions}
\label{sec:res-ml-int-research-questions}

% TODO: State the questions that must be answered before the Deliverable
% can be accepted.

\subsection{Terminology and Definitions}
\label{sec:res-ml-int-definitions}

% TODO: Define the technical terminology, symbols, quantities and
% classifications used in this Deliverable.

\subsection{Literature Evidence}
\label{sec:res-ml-int-literature-evidence}

% TODO: Record evidence from peer-reviewed papers, authoritative datasets,
% standards, technical reports and primary documentation.
%
% Do not create a detached literature-summary list. Explain how each source
% supports, contradicts or limits a project decision.

\subsection{Governing Relations and Quantitative Definitions}
\label{sec:res-ml-int-governing-relations}

% TODO: Record relevant equations, dimensionless groups, empirical models,
% extraction rules or quantitative definitions.
%
% Use align environments for displayed equations.
% Every displayed equation must have a unique \label{eq:...}.
% Do not place punctuation after displayed equations.

\subsection{Machine-Learning Role}
\label{sec:res-ml-int-machine-learning-role}

% TODO: Complete this RES-ML Domain-specific subsection.

\subsection{Non-ML Baseline}
\label{sec:res-ml-int-non-ml-baseline}

% TODO: Complete this RES-ML Domain-specific subsection.

\subsection{Input and Output Representation}
\label{sec:res-ml-int-input-output-representation}

% TODO: Complete this RES-ML Domain-specific subsection.

\subsection{Dataset Requirements}
\label{sec:res-ml-int-dataset-requirements}

% TODO: Complete this RES-ML Domain-specific subsection.

\subsection{Model or Architecture Candidates}
\label{sec:res-ml-int-model-architecture-candidates}

% TODO: Complete this RES-ML Domain-specific subsection.

\subsection{Training and Validation Method}
\label{sec:res-ml-int-training-validation-method}

% TODO: Complete this RES-ML Domain-specific subsection.

\subsection{Evaluation Metrics}
\label{sec:res-ml-int-evaluation-metrics}

% TODO: Complete this RES-ML Domain-specific subsection.

\subsection{Generalisation and Uncertainty}
\label{sec:res-ml-int-generalisation-uncertainty}

% TODO: Complete this RES-ML Domain-specific subsection.

\subsection{Simulation and Optimisation Integration}
\label{sec:res-ml-int-simulation-optimisation-integration}

The deferred integration sequence is:
\begin{enumerate}
  \item consume validated Tohoku corridor and barrier-class outputs;
  \item specify a future candidate configuration or intervention;
  \item verify that the query lies inside the future surrogate domain;
  \item predict response and uncertainty;
  \item screen or rank reinforcement options;
  \item require direct hydrodynamic, structural or FSI simulation when
    uncertain or out of domain;
  \item add accepted verified simulations to the controlled dataset.
\end{enumerate}

\subsection{Candidate Approaches}
\label{sec:res-ml-int-candidate-approaches}

% TODO: Define the credible candidate approaches considered.

\subsection{Critical Comparison}
\label{sec:res-ml-int-critical-comparison}

% TODO: Compare candidates using physically and project-relevant criteria.
% Do not select an approach solely on popularity or implementation ease.

\subsection{Assumptions and Applicability}
\label{sec:res-ml-int-assumptions}

% TODO: State assumptions and the conditions under which they are valid.

\subsection{Limitations and Failure Conditions}
\label{sec:res-ml-int-limitations}

% TODO: State where the evidence, model or selected approach ceases to be
% reliable or applicable.

\subsection{Project-Specific Conclusions}
\label{sec:res-ml-int-conclusions}

Integration remains a governed future screening workflow. For the
current phase, accepted physical evidence flows from validated
simulation into structured outputs only; surrogate predictions do not
route decisions.

\subsection{Selected Approach and Rationale}
\label{sec:res-ml-int-selected-approach}

% TODO: Record the selected approach only after sufficient evidence exists.
% State the alternatives rejected and the reasons for rejection.

\subsection{Unresolved Decisions}
\label{sec:res-ml-int-unresolved-decisions}

% TODO: Record unresolved questions, required evidence and decision owners.

\subsection{Requirements and Handoffs}
\label{sec:res-ml-int-requirements}

% TODO: State requirements passed to other Research Domains, Software,
% verification activities or later execution work.

\subsection{Deliverable Completion Record}
\label{sec:res-ml-int-completion-record}

% TODO: Record whether the Deliverable is:
%
% - Not started
% - In progress
% - Draft complete
% - Reviewed
% - Accepted
%
% Acceptance requires:
%
% - the research questions to be answered;
% - claims to be supported by appropriate evidence;
% - assumptions and limitations to be explicit;
% - the project-specific conclusion to be stated;
% - downstream requirements to be traceable;
% - unresolved decisions to be closed or formally deferred.

% -------------------------------------------------------------------------
% WBS ID: RES-ML-SPEC
% Deliverable: Integrated Machine-Learning Specification
% Status: In progress
% Evidence status: ML deferred by current hydrodynamic scope
% Review status: Not reviewed
% Last updated: 2026-07-19
% Dependencies: RES-ML-ROL through RES-ML-INT, RES-DAT-CAS, RES-GEO-SPEC, RES-OPT-SPEC
% Downstream interfaces: GitHub research review, deferred Jira WBS update
% -------------------------------------------------------------------------

\section{RES-ML-SPEC --- Integrated Machine-Learning Specification}
\label{sec:res-ml-spec}

\subsection{Purpose and Scope}
\label{sec:res-ml-spec-purpose}

This Deliverable records that the event-conditioned defence-assessment
ML specification is deferred under the current hydrodynamic scope. It
preserves the WBS ID without creating a replacement hierarchy.

\subsection{Research Questions}
\label{sec:res-ml-spec-research-questions}

% TODO: State the questions that must be answered before the Deliverable
% can be accepted.

\subsection{Terminology and Definitions}
\label{sec:res-ml-spec-definitions}

% TODO: Define the technical terminology, symbols, quantities and
% classifications used in this Deliverable.

\subsection{Literature Evidence}
\label{sec:res-ml-spec-literature-evidence}

% TODO: Record evidence from peer-reviewed papers, authoritative datasets,
% standards, technical reports and primary documentation.
%
% Do not create a detached literature-summary list. Explain how each source
% supports, contradicts or limits a project decision.

\subsection{Governing Relations and Quantitative Definitions}
\label{sec:res-ml-spec-governing-relations}

% TODO: Record relevant equations, dimensionless groups, empirical models,
% extraction rules or quantitative definitions.
%
% Use align environments for displayed equations.
% Every displayed equation must have a unique \label{eq:...}.
% Do not place punctuation after displayed equations.

\subsection{Machine-Learning Role}
\label{sec:res-ml-spec-machine-learning-role}

Any future ML specification shall consolidate candidate ML roles,
historical evidence, simulation-derived data, barrier and intervention
representations, algorithm hierarchy, losses and labels, grouped
validation, uncertainty, out-of-domain handling, simulation fallback,
reinforcement-screening workflow, limitations and extension gates.

\subsection{Non-ML Baseline}
\label{sec:res-ml-spec-non-ml-baseline}

% TODO: Complete this RES-ML Domain-specific subsection.

\subsection{Input and Output Representation}
\label{sec:res-ml-spec-input-output-representation}

% TODO: Complete this RES-ML Domain-specific subsection.

\subsection{Dataset Requirements}
\label{sec:res-ml-spec-dataset-requirements}

% TODO: Complete this RES-ML Domain-specific subsection.

\subsection{Model or Architecture Candidates}
\label{sec:res-ml-spec-model-architecture-candidates}

% TODO: Complete this RES-ML Domain-specific subsection.

\subsection{Training and Validation Method}
\label{sec:res-ml-spec-training-validation-method}

% TODO: Complete this RES-ML Domain-specific subsection.

\subsection{Evaluation Metrics}
\label{sec:res-ml-spec-evaluation-metrics}

% TODO: Complete this RES-ML Domain-specific subsection.

\subsection{Generalisation and Uncertainty}
\label{sec:res-ml-spec-generalisation-uncertainty}

% TODO: Complete this RES-ML Domain-specific subsection.

\subsection{Simulation and Optimisation Integration}
\label{sec:res-ml-spec-simulation-optimisation-integration}

% TODO: Complete this RES-ML Domain-specific subsection.

\subsection{Candidate Approaches}
\label{sec:res-ml-spec-candidate-approaches}

% TODO: Define the credible candidate approaches considered.

\subsection{Critical Comparison}
\label{sec:res-ml-spec-critical-comparison}

% TODO: Compare candidates using physically and project-relevant criteria.
% Do not select an approach solely on popularity or implementation ease.

\subsection{Assumptions and Applicability}
\label{sec:res-ml-spec-assumptions}

% TODO: State assumptions and the conditions under which they are valid.

\subsection{Limitations and Failure Conditions}
\label{sec:res-ml-spec-limitations}

% TODO: State where the evidence, model or selected approach ceases to be
% reliable or applicable.

\subsection{Project-Specific Conclusions}
\label{sec:res-ml-spec-conclusions}

% TODO: Convert the evidence into conclusions specific to the Studentship
% project. Do not merely restate literature findings.

\subsection{Selected Approach and Rationale}
\label{sec:res-ml-spec-selected-approach}

The selected active approach is no ML implementation or training.
Future event-conditioned supervised learning may be reconsidered only
after the Tohoku hydrodynamic workflow, fixed-rigid barrier metrics and
V1--V2--V3 validation evidence are accepted.

\subsection{Unresolved Decisions}
\label{sec:res-ml-spec-unresolved-decisions}

% TODO: Record unresolved questions, required evidence and decision owners.

\subsection{Requirements and Handoffs}
\label{sec:res-ml-spec-requirements}

The integrated specification passes the following requirements downstream:
\begin{itemize}
  \item no ML training before event, local hydrodynamic and
    impact/loading validation gates;
  \item no claim that historical geometry changes are optimal causal labels;
  \item grouped validation by complete simulation case;
  \item explicit uncertainty and out-of-domain rejection;
  \item direct hydrodynamic, structural or FSI fallback for uncertain
    or unsupported future queries;
  \item deferred Jira WBS updates through the pending change manifest.
\end{itemize}

\subsection{Deliverable Completion Record}
\label{sec:res-ml-spec-completion-record}

% TODO: Record whether the Deliverable is:
%
% - Not started
% - In progress
% - Draft complete
% - Reviewed
% - Accepted
%
% Acceptance requires:
%
% - the research questions to be answered;
% - claims to be supported by appropriate evidence;
% - assumptions and limitations to be explicit;
% - the project-specific conclusion to be stated;
% - downstream requirements to be traceable;
% - unresolved decisions to be closed or formally deferred.


\clearpage
\printbibliography[
  heading=bibintoc,
  title={References}
]

\end{document}
